\documentclass[../main.tex]{subfiles}
\begin{document}
%==========================================TIÊU ĐỀ TẬP========================================
\begin{table}[h]
  \begin{adjustwidth}{-1cm}{}
    \begin{tabular}{>{\centering\arraybackslash}p{6cm}|>{\raggedright\arraybackslash}p{12.5cm}}
      \multirow{5}{6cm}
      % Cột bên trái: ÉPISODE và số tập
      {\\[-40pt]
        \flaregothic\fontsize{20pt}{20pt}\selectfont \centering
        \color{eptype}ÉPISODE\color{black}\\
        \burbank\fontsize{72pt}{72pt}\selectfont
        \color{epnum}46\color{black}\\[6pt]
        \cafeteria\fontsize{20pt}{20pt}\selectfont\color{epnum}(4-4)\color{black}
        \\[18pt]
        \flaregothic\fontsize{18pt}{18pt}\selectfont
        \color{epattr}ÉPISODE PERDU\color{black}
      }
       &
      % Dòng thứ nhất: tên tập tiếng Nhật
      {\fotskip\fontsize{18pt}{18pt}\selectfont パワハラの\ruby{現}{げん}\ruby{場}{ば}をおさえました
      }                                            
                                  \\[6pt]
      % Dòng thứ hai: tên tập tiếng Anh
       & {\cafeteria\fontsize{18pt}{18pt}\selectfont Factory Work is a Hard Work}                                        
                                       \\[4pt]
      % Dòng thứ ba: tên tập tiếng Việt
       & {\vnmsans\fontsize{18pt}{18pt}\selectfont Lao động (trong game) là vinh quang!}                                  
                                      \\[8pt]
      % Dòng thứ tư: Nhân vật xuất hiện
       & {\caviar{\fontsize{14pt}{14pt}\selectfont Nhân vật: \emp{kuon1} \emp{ryuuhou1} \emp{mel} \emp{choco} \emp{subaru} \emp{pekora} \emp{marine}}}
      \\[8pt]
      % Dòng cuối cùng: Thời gian diễn ra của tập (thường sẽ là ngày phát hành)
       & {\continuum\fontsize{16pt}{16pt}\selectfont Ngày phát hành: 22/03/2020}                                                                               \\[18pt]
    \end{tabular}
  \end{adjustwidth}
\end{table}
%=========================================TRUYỆN CHÍNH========================================
\color{black}

\txtbx[2]{{\color{vol} \uline{Tóm tắt:}} Căn bệnh nghiện game của Pekora bất ngờ tái phát với một chiếc máy hỏng, vô tình hút Subaru, Mel và Marine vào một thế giới RPG kỳ lạ. Tại đây, thay vì làm anh hùng cứu thế giới, họ lại bị bóc lột sức lao động trong một nhà máy sản xuất sữa lắc dưới trướng giám sát viên Marine. Kuon đành phải cầu cứu cô em gái game thủ Ryuuhou đến ``gánh tạ'' phá đảo trò chơi, từ đó mới bộc lộ ra tính cách cộc cằn của cô bé.}

Một ngày nọ, văn phòng \hll\ lại phải đối mặt với một thảm họa quen thuộc: Pekora lại tái phát căn bệnh nghiện trò chơi điện tử. 

Cô nàng Thỏ xanh ngồi bệt xuống mặt sàn trước chiếc tivi lớn, hai tay nắm chặt rịt lấy chiếc tay cầm của máy console đã bị Rushia lỡ tay đập hỏng từ hôm trước. Dẫu cho dây điện đã đứt lìa và chiếc máy rõ ràng không còn khả năng hoạt động, Pekora vẫn không ngừng bấm các nút bấm kêu tanh tách liên hồi. Ánh mắt cô vô hồn, mờ mịt, hai chiếc tai thỏ rũ cúp xuống một cách mệt mỏi, trông như thể toàn bộ thế giới xung quanh cô lúc này chỉ còn thu bé lại vừa bằng một màn hình trò chơi tưởng tượng.

\img{4-4a}

Subaru vừa đẩy cửa bước vào phòng nghỉ, liền trợn tròn mắt kinh ngạc trước cảnh tượng kỳ dị: ``Ê này, Pekora-chan đang ngồi bấm khí thế trên cái máy bị hỏng tanh bành hôm bữa kìa! Mà máy hỏng rồi thì cắm điện kiểu gì mà chơi được nhỉ?''

Đi ngay phía sau lưng nàng Tomboy, Mel và Marine cũng đưa mắt nhìn cô nàng Thỏ với vẻ cực kỳ khó hiểu. Họ hoàn toàn không rõ chuyện quái quỷ gì đang xảy ra, chỉ thấy sự lạ lùng toát ra từ biểu cảm vô hồn như cá chết của Pekora. 

Ngay lúc đó, Kuon từ ngoài hành lang bước vào, trên tay cậu vẫn còn cầm kẹp tài liệu sau buổi họp giao ban. Vừa nhìn thấy cảnh tượng đập vào mắt, cậu liền lắc đầu ngao ngán, đưa tay lên bóp trán: ``Thôi xong, lại nữa rồi\dots''

Thuyền trưởng thấy sự ngán ngẩm hiện rõ mồn một trên khuôn mặt của vị Tổng quản lý, liền tò mò hỏi: ``Có chuyện gì xảy ra với cậu ấy sao, Kuon-senpai?''

Cậu nhún vai, ánh mắt toát lên sự bất lực tột độ: ``Anh chịu chết. Lúc anh vừa mới đến đây khoảng nửa giờ trước là em ấy đã ngồi ru rú ở cái góc đó với tình trạng này rồi. Chắc là căn bệnh nghiện game từ đợt trước lại tái phát nặng hơn rồi đây.''

``Ý anh nói tái phát là sao cơ?'' Subaru nhướng mày thắc mắc.

Cậu quản lý thở dài nặng nề một hơi, đặt cặp tài liệu cộp xuống bàn rồi quay sang nhìn cả nhóm: ``Thôi được rồi, để anh kể tóm tắt lại. Nhưng cấm có đứa nào được phàn nàn là anh nói dài dòng đấy nhé\dots''

Ba cô gái kéo ghế ngồi xung quanh, chăm chú lắng nghe Kuon kể lại chuỗi sự kiện hỗn loạn xảy ra vài tuần trước. Khi đó, Shion, Choco, Okayu, Rushia và chính cậu đã phải bày đủ mọi chiêu trò để giúp Pekora ``chữa'' căn bệnh nghiện game\footnote{Xem lại {\flaregothic ÉPISODE 36}, quyển số Ba.}.

``\dots và mọi chuyện là như vậy đấy. Dù bọn anh đã cố gắng ép em ấy chơi những trò chơi vận động thể chất ngoài đời thực không đụng đến đồ điện tử, nhưng rốt cuộc em ấy vẫn ngựa quen đường cũ, tái nghiện như bình thường,'' cậu kết thúc dòng hồi tưởng đầy mệt mỏi của mình. ``Có vẻ như đợt trước bọn anh mới chỉ cắt được phần ngọn, chứ chưa thể nhổ tận gốc cái mầm mống nghiện ngập này.''

Sau khi nghe xong lời bộc bạch từ cậu, Mel khẽ gật đầu, rồi liếc đôi mắt thương cảm về phía Pekora đang ngồi run lẩy bẩy góc phòng: ``Ra là vậy\dots\ Em ấy đúng là nghiện nặng đến mức hết thuốc chữa rồi\dots''

``Chả trách mà nãy giờ em thấy cậu ấy cứ ngồi đờ đẫn ra như người mất hồn,'' Subaru chêm vào, giọng điệu tỏ vẻ vừa tội nghiệp vừa bất lực với cô bạn đồng nghiệp.

Marine đưa tay chống cằm, làm ra vẻ suy nghĩ vô cùng sâu xa, rồi ngước lên nhìn Kuon: ``Vậy, Kuon-senpai này, lần này chúng ta nên triển khai phương án gì đây? Anh có cao kiến nào không?''

Cậu Tổng liếc nhìn nàng Thỏ xanh đang bấm nút trong vô thức, rồi khoanh hai tay lại trước ngực, nhắm nghiền mắt phân tích: ``Anh nghĩ là lần này chúng ta phải tìm ra một cách tiếp cận hoàn toàn mới thôi. Nhưng muốn làm gì thì làm, nhiệm vụ tối thượng bây giờ là phải lôi xềnh xệch được em ấy ra khỏi cái máy chơi game ảo ma kia đã.''

Cả bốn người xúm lại bắt đầu bàn bạc rôm rả, ai nấy đều hăng hái đưa ra những ý tưởng táo bạo của riêng mình để kéo cô nàng Thỏ trở lại thực tại. Nhưng cứ mỗi lần nhìn sang cái cách Pekora vẫn đang ngồi bần thần, tay bấm nhầm nút liên tục trên chiếc tay cầm đứt dây, họ đều tự nhận thức được rằng nhiệm vụ này hoàn toàn không hề dễ xơi chút nào.

\ct

Trong lúc cả bốn người họ vẫn còn đang tranh luận nảy lửa, Mel khẽ nghiêng đầu, đôi mắt màu vàng hổ phách toát lên vẻ tò mò khi quan sát hành động của Pekora. 

``Nói về mấy cái tựa trò chơi trên máy console thì\dots\ chắc là em ấy đang tưởng tượng mình đang chơi mấy trò nhập vai (RPG) chăng? Cái kiểu nhập vai thành một vị anh hùng dũng cảm cầm kiếm đi cứu thế giới ấy?'' nàng Dơi đưa ra giả thuyết.

Kuon đang đứng khoanh tay dựa lưng vào tường, khẽ gật gù đồng ý: ``Cũng dễ là thế lắm. Tiếc là mấy cái thể loại chặt chém cày cuốc này anh lại không rành cho lắm. Nhưng mà theo anh để ý, mấy cái motif trò chơi nhập vai này thường thì người chơi toàn được làm anh hùng xuất chúng đi cứu công chúa, đúng không?''

Subaru và Marine nghe cậu con trai nhắc đến hai chữ ``anh hùng'' mà bất giác mỉm cười đầy mơ mộng. Trong thâm tâm hai cô nàng bỗng dấy lên một thứ cảm giác hoài niệm mãnh liệt về những ngày tháng tuổi thơ vô tư, khi họ vẫn còn ôm ấp những giấc mộng lớn lao.

``Hầy\dots\ Tự dưng nhắc lại làm em nhớ những ngày mình còn nhỏ, lúc nào cũng ước ao muốn được trở thành một anh hùng quá đi mất,'' nàng Tomboy thở dài, ánh mắt xa xăm hoài niệm lại những kỷ niệm xưa cũ.

Nàng Hải tặc cũng cười mỉm, gật đầu cái rụp đồng tình: ``Đám con gái chúng mình, ai mà chẳng có mộng mơ muốn được làm anh hùng vĩ đại một lần trong đời chứ lị. Hồi bé em toàn tưởng tượng mình là một nữ chiến binh ngầu lòi, vác kiếm đi bảo vệ nền hòa bình của thế giới cơ đấy.''

\img{4-4b}

Nghe hai người em say sưa kể chuyện, nàng Dơi liền quay sang cậu Tổng, vẻ mặt đầy hiếu kỳ: ``Thế còn Kuon-senpai thì sao ạ? Hồi bé anh có hay mơ mộng được làm siêu anh hùng giải cứu thế giới không?''

Kuon khẽ nhướng một bên mày, rồi chậm rãi thở ra một hơi dài, như thể đang moi móc lại một miền ký ức vô cùng xa xôi. ``Ờm\dots\ hoàn toàn không. Ngày trước anh toàn cắm mặt vào mấy cuốn sách giải đố logic là nhiều. Với lại gia cảnh nhà anh lúc đó cũng chẳng khá giả gì để sắm cái máy chơi game console nào cả, nên anh cũng chẳng có cơ hội mà cày game mơ mộng như mấy đứa đâu. Ước mơ lớn nhất hồi đó của anh chỉ đơn giản là giải mã được hết mấy cái câu đố hóc búa trong cuốn bách khoa toàn thư thôi.''

Subaru quay ngoắt sang nhìn cậu, vẻ mặt vừa ngạc nhiên tột độ vừa có chút buồn cười: ``Cái gì cơ chứ!? Vậy mà xưa nay em cứ tưởng anh phải là cái kiểu người đam mê sức mạnh, thích làm siêu anh hùng trừ gian diệt ác cơ đấy!''

Cậu Tổng nhún vai, nhếch mép cười nhạt: ``Anh sống theo chủ nghĩa thực tế mà em, làm siêu anh hùng đi gánh vác thế giới chỉ tổ rước họa vào thân, mệt mỏi thêm thôi. Làm một người bình thường sáng cắp ô đi tối cắp ô về giải đố mua vui có phải nhàn nhã hơn nhiều không.''

Marine bật cười khúc khích trước tính cách ``cụ non'' của cậu, còn Mel thì nhìn Kuon bằng một ánh mắt vô cùng cảm thông: ``Vậy chắc hẳn hồi bé tuổi thơ của anh ít được trải nghiệm mấy trò phiêu lưu mạo hiểm lắm nhỉ?''

``Ừ. Nhưng mà bù lại thì bây giờ anh lại đang phải trải nghiệm cảm giác phiêu lưu sinh tồn ngoài đời thực với cái văn phòng toàn mấy đứa siêu quậy đây này,'' cậu thở dài não nề, ngẫm lại những chuỗi sự kiện phá hoại banh chành công ty đã qua.

Lời cảm thán đầy chua xót của Kuon khiến cả ba cô gái bật cười nắc nẻ, làm cho bầu không khí căng thẳng trong phòng có phần nhẹ nhàng và dễ thở hơn. Thế nhưng, ánh mắt của mọi người vẫn không quên đổ dồn về phía Pekora đang hoàn toàn đắm chìm vào trong thế giới ảo giác của riêng mình. Họ biết rõ rằng để kéo được cô nàng Thỏ ương bướng này ra khỏi tình trạng ảo ma này sẽ cần phải có một kế hoạch tác chiến cực kỳ bài bản.

Thế nhưng, ngay khi mọi người vừa dứt tiếng cười, cái máy chơi game tàn tạ, đứt dây nằm trên kệ tivi đó bỗng nhiên phát ra một giọng nói điện tử ồm ồm rợn tóc gáy: ``Nếu đã như vậy\dots\ ta sẽ biến điều ước sâu thẳm của các ngươi thành hiện thực.''

Cả bốn người nghe xong liền giật bắn mình, đứng sững lại như trời trồng khi nghe thấy âm thanh phát ra từ một thiết bị tưởng chừng như đã chết hoàn toàn.

Subaru hoảng hốt lùi lại phía sau, lắp bắp thốt lên: ``Oái! Cái\dots\ cái máy phế liệu này nó biết nói tiếng người hả!?''

Mel run rẩy núp sau lưng Marine: ``Chuyện quái quỷ gì đang xảy ra thế này mọi người ơi?''

Marine thì bấu chặt lấy cánh tay của cậu Tổng, giọng run lên bần bật: ``N-Này, Kuon-senpai\dots\ anh đừng có bảo với em là cái máy này\dots\ là hàng bị dính nguyền rủa tâm linh nhé!?''

Kuon vội gạt tay cô nàng ra, đổ mồ hôi hột đáp: ``Anh cũng có biết quái đâu! Nhưng mà cái tình huống tâm linh này đúng là vượt quá sức tưởng tượng của con người rồi đấy\dots''

Vừa dứt câu, một luồng ánh sáng trắng chói lóa phát ra từ khe cắm đĩa của chiếc máy, nhanh chóng bao trùm nuốt chửng lấy toàn bộ căn phòng. Giữa luồng sáng chói lóa ấy, Pekora vẫn ngồi im thít cầm chiếc tay cầm, một nụ cười ngờ nghệch, vô tri vẫn hiện rõ trên khuôn mặt với đôi mắt trống rỗng. 

Đến khi ánh sáng ma thuật vụt tắt hẳn, căn phòng nghỉ giờ đây chỉ còn trơ lại mỗi cậu Tổng quản lý cùng nàng Thỏ, và chiếc máy chơi điện tử thì đã thần kỳ trở lại nguyên trạng bóng loáng hệt như chưa từng bị hỏng hóc gì.

\begin{figure}[H]
  \centering
  \begin{subfigure}[t]{0.45\textwidth}
    \centering
    \includegraphics[width=8cm]{../../../../Image/Book4/4-4c1.png}
  \end{subfigure}
  \quad
  \begin{subfigure}[t]{0.45\textwidth}
    \centering
    \includegraphics[width=8cm]{../../../../Image/Book4/4-4c2.png}
  \end{subfigure}
\end{figure}

Tuy nhiên, khi ngước mắt nhìn lên màn hình tivi, Kuon lập tức phát hiện ra một điều khiến cậu toát mồ hôi lạnh sống lưng. Một tựa trò chơi nhập vai mới toanh vừa xuất hiện trên màn hình hiển thị dòng chữ tiêu đề:

\begin{center}
  \uline{``Ba nàng lính\dots\ mộng mơ\dots?''}
\end{center}

Cậu giật thót mình lùi lại vài bước, hoảng hốt đưa mắt nhìn quanh quất khắp phòng nhưng tuyệt nhiên không hề nhìn thấy bóng dáng của ba cô gái kia đâu cả. Cậu há hốc miệng lẩm bẩm: ``\textit{Đồng cây mát mẻ\dots\ Đừng có nói với mình là Mel, Marine với Subaru vừa bị hú vào trong thế giới game ảo rồi đấy nhé!?}''

Pekora vẫn ngồi ru rú ở đó, hai tay ôm khư khư lấy chiếc tay cầm với vẻ mặt vô hồn. Kuon vội vã quay phắt sang túm vai nàng Thỏ lay mạnh, nhưng trước khi cậu kịp mở miệng chất vấn, cô nàng đã lẩm bẩm trong vô thức: ``Peko peko\dots\ cứu họ đi, Kuon-senpai\dots'' Dường như sâu thẳm bên trong tiềm thức, cô nàng vẫn còn giữ lại được một chút tỉnh táo và có thể mường tượng được cái tai họa tày đình vừa mới xảy ra.

Cậu Tổng buông thõng tay, chán nản vò đầu bứt tóc xoa trán đầy bất lực. ``Mà khổ nỗi cái thân mình xưa nay có giỏi ba cái thể loại RPG chặt chém quái quỷ này đâu cơ chứ\dots\ Nếu muốn phá đảo game cứu người thì chắc phải gọi ngay cho đám GAMERS nhờ vả mới xong\dots'' Cậu vội vã thọc tay vào túi quần rút điện thoại ra định gọi cứu viện, nhưng rồi lại sực nhớ ra một sự thật phũ phàng: ``Chết dở, giờ này cả bốn người họ lại đang bận bịu phát sóng livestream hết rồi mới cay chứ! Giờ biết gọi ai vào gánh tạ bây giờ\dots''

Tưởng chừng như một mình cậu Tổng dân IT sẽ phải ngậm đắng nuốt cay tự mình cày cuốc với tựa trò chơi ``trái ngành'' này, thì đột nhiên, ánh mắt cậu sáng rực lên như vừa bắt được phao cứu sinh.  ``A! Phải rồi! Ryuuhou! Gọi ngay em ấy vác mặt đến đây gánh team là chuẩn bài nhất! Nó giỏi mấy cái kiểu này mà!''

Lướt tay nhanh thoăn thoắt qua danh bạ điện thoại, cậu lập tức bấm số gọi cho cô em gái ruột của mình - một game thủ hạng nặng kiêm trưởng nhóm nhạc underground. Chưa đầy ba giây sau, một cuộc gọi video đã được kết nối, để lộ ra khuôn mặt đang nhăn nhó cọc cằn cùng mái tóc rối bời của cô em gái. Giọng con bé vang lên đầy vẻ bực dọc, gắt gỏng: ``Gọi gì vào cái giờ thiêng này hả ông anh!? Có biết là em đang dở trận chuỗi lên rank không!? Có rớt hạng thì anh đền nổi không hả?''

\chrint

Nếu như trong mắt mọi người, Hikawa Ryuuhou thường ngày luôn được biết đến là một cô bé tomboy hướng ngoại, năng động, hay cười và thỉnh thoảng có những trò đùa nhí nhảnh, thì khi bước vào thế giới của các trò chơi điện tử, cô bé lại biến thành một con người hoàn toàn khác, không khác nào có một tính cách thứ hai bên trong, ngay cả với cách xưng hô.

Người ta thường nói chơi game là để giải trí, nhưng đối với Ryuuhou, game là một chiến trường sinh tử mang đậm tính sát phạt. Khát vọng chiến thắng và máu ăn thua của cô bé cao đến mức cực đoan. Chỉ cần ngón tay chạm vào chiếc tay cầm, bàn phím hay điện thoại thôi, nụ cười hiền lành thường ngày lập tức bị xóa sổ, thay vào đó là ánh mắt sắc lẹm hình viên đạn, đôi lông mày nhíu chặt lại và tư thế ngồi chồm hẳn về phía trước đầy sát khí.

Trong game, Ryuuhou là một ``kẻ độc tài'' chính hiệu. Cô bé chơi game với một sự tập trung cao độ đến đáng sợ, và đi kèm với đó là cái nết vô cùng cọc cằn, dễ bốc hỏa. Chỉ cần đồng đội di chuyển sai một nhịp, hệ thống game bị giật lag, hay kịch bản game có phần vô lý, Ryuuhou sẵn sàng buông ra những lời chửi đổng, văng tục ``giảm xóc'' liên thanh không ngừng nghỉ. Khi chiến thắng, con bé sẽ cười ngạo nghễ và ``gáy'' cực kỳ khét; nhưng nếu lỡ thua hoặc bị game làm cho ức chế, hậu quả nhẹ thì đập bàn đập ghế, nặng thì những chiếc tay cầm hoặc bàn phím đó xác định sẽ biến thành đống sắt vụn dưới cơn thịnh nộ của nàng game thủ. 

Nhưng may mắn thay, tính cách đó không kéo dài quá lâu. Mỗi khi kết thúc một trận game căng thẳng, dù thắng hay thua, Ryuuhou đều có thói quen hít một hơi thật sâu, bẻ khớp tay răng rắc, rồi dùng hai tay vỗ mạnh vào má mình. Cái tiếng *BỘP* vang lên đó chính là âm thanh ``khởi động lại'' tính cách nhí nhảnh của cô bé. Chỉ cần sau tiếng vỗ đó, Ryuuhou quay sang nhìn Kuon và hỏi: ``Onii-san, nay nhà mình ăn gì?``, thì đó là lúc ``nhân cách game thủ'' đã được tắt thành công.

Chính vì cái ``nhân cách thứ hai'' trái ngược hoàn toàn này mà Kuon thường xuyên phải đóng vai trò làm bảo mẫu dọn dẹp bãi chiến trường, hoặc phải chật vật làm dịu đi cái đầu đang bốc hỏa của cô em gái mỗi khi con bé lỡ chơi quá đà. Kuon vẫn thường hay đùa với mọi người ở công ty rằng: ``''\textit{Nếu muốn biết bản chất thật sự của một đứa con gái, hãy nhìn cái cách mà nó chơi game khi mạng bị rớt ping.}''

\chrint

Kuon vội vàng gạt đi, không để cho cô em có cơ hội cằn nhằn thêm: ``Ryuuhou! Bỏ ngay cái trận rank rác rưởi đó đi! Anh đang có việc khẩn cấp liên quan đến tính mạng cần em vác mặt đến đây giúp ngay bây giờ!''

Thấy điệu bộ nghiêm trọng hốt hoảng của ông anh trai, Ryuuhou cũng đành tạm gác lại cơn bực, nhíu mày hỏi: ``Được rồi, được rồi, ông anh bình tĩnh xem nào. Tự dưng làm quá lên. Có chuyện gì mà nghe như cháy nhà thế?''

``Mấy đứa nhân viên của công ty anh vừa bị cái máy chơi game nguyền rủa hút vào trong thế giới ảo rồi! Em là cao thủ RPG, anh cần em đến đây gánh anh phá đảo trò này ngay lập tức để lôi mấy đứa ấy ra ngoài!'' Kuon nói một tràng dài với vẻ mặt cắt không còn giọt máu.

Nghe lời giải thích nghe đậm mùi viễn tưởng của ông anh, Ryuuhou ngớ người ra, nhìn chằm chằm vào màn hình điện thoại với biểu cảm như đang nhìn một thằng dở hơi. ``Ông anh đang bú đá ảo giác hay gì đấy? Nói cái méo gì nghe viễn tưởng vãi nho thế?''

``Anh mày rảnh méo đâu mà đùa giỡn với mày giờ này! Chuyện tâm linh dài dòng lắm, anh sẽ kể chi tiết cho em nghe sau!'' Kuon hối thúc liên tục.

Cảm nhận được sự khẩn trương có một không hai của ông anh Tổng quản lý, cô nàng game thủ cá tính cũng đành phải thở dài nhượng bộ, cố gắng tỏ vẻ điềm tĩnh chuyên nghiệp: ``Thôi được rồi. Thế cái tựa game ma cỏ đó tên là gì? Cứ mở máy bật màn hình lên trước đi. Em sẽ nán lại nhà cắm mic hướng dẫn anh chơi từ xa\dots''

Thế nhưng, Kuon lập tức cau mày khi nhìn xuống dòng chữ cảnh báo bé tí in ngay dưới tên của trò chơi trên màn hình: ``Không được đâu em ơi\dots\ Lần này anh bắt buộc phải cần em lết xác đến tận văn phòng. Cái con game hãm này\dots\ nó yêu cầu phải có đủ ba người cùng chơi (Co-op 3 players) thì mới cho vào game. Hiện tại thì Peko-chan gây họa đang ngồi thẫn thờ ở đây giữ một chỗ rồi, anh không thể nào tự solo một mình hai cái tay cầm được.''

``Bắt buộc ba người chơi á? Cái thể loại game rác rưởi vắt sữa kiểu gì mà thiết kế cái luật ngu học thế?'' Ryuuhou chửi thề một tiếng.

``Anh cũng chả biết được, nhưng nhìn vào cái giao diện sảnh chờ thì anh đoán là mỗi người sẽ phải ôm một cái tay cầm để điều khiển trực tiếp một nhân vật tương ứng. Nếu thiếu dù chỉ một người, trò chơi sẽ không cho phép ấn nút Start. Em đến đây lẹ lẹ giúp anh, nếu không thì ba đứa kia sẽ vĩnh viễn bốc hơi kẹt luôn trong cái thế giới ảo đó mất!''

Ryuuhou khẽ tặc lưỡi, nhếch mép cười một cái đầy hiếu chiến: ``Ồ, nghe mùi cũng kịch tính và ảo ma đấy. Được rồi, nhắn ngay định vị cái văn phòng mới qua đây. Bà mày sẽ xách xe chạy qua gánh tạ cho ông anh ngay bây giờ.'' 

Sau khi chốt hạ được đối phương đồng ý tới cứu viện, Kuon mới dám thở phào nhẹ nhõm vuốt mồ hôi trán: ``Cứu tinh đây rồi! Nhanh cái chân lên đấy!''

\ct

\begin{center}
  \uline{(Phần nội dung dưới đây diễn biến bên trong tựa game.\\Kuon, Ryuuhou và Pekora đóng vai trò là những bình luận viên ngồi trước màn hình.)}
\end{center}

\img{4-4d}

Khung cảnh đầu tiên hiện ra là một thế giới hoang tàn bị thế lực tà ác tàn phá nặng nề, nơi bóng tối mịt mù bao trùm che lấp khắp các vùng đất, gieo rắc sự chết chóc và nỗi tuyệt vọng cùng cực lên nhân loại. Giữa đống tro tàn đó, một giọng đọc truyền thuyết hào hùng cất lên kể rằng chỉ có duy nhất một vị anh hùng dũng cảm được số phận lựa chọn mới có đủ khả năng tiêu diệt thế lực bóng tối này và mang lại nền hòa bình vĩnh cửu.

Ngay sau đó, một luồng hào quang xuất hiện, ánh sáng thần thánh chói lóa lấp đầy toàn bộ màn hình, và một dòng chữ hoành tráng được chạy dọc hiện lên: ``\uline{Vị anh hùng Subaru, kẻ được định mệnh vinh danh lựa chọn để gánh vác sứ mệnh cứu thế giới!}''

\il

Ryuuhou vừa nhai kẹo cao su chóp chép vừa cười khẩy mỉa mai: ``Hờ, cái mở bài dẫn truyện làm màu nghe oai phết nhỉ. Để xem con nhỏ Vịt này sẽ được thiết kế xuất hiện ngầu lòi thế nào! Chắc lại kiểu từ trên trời giáng xuống cầm thanh gươm ánh sáng chặt chém ra trò đây.''

Kuon gật gù phụ họa: ``Thì motif anh hùng nào lúc mở bát chả phải làm màu để loe mắt thiên hạ.''

\il

Thế nhưng, cảnh quay tiếp theo hiện ra lại tát thẳng một gáo nước lạnh vào mặt người xem. Subaru, lúc này đang khoác trên mình bộ đồng phục công nhân nhà máy xanh lè với chiếc mũ bảo hộ màu trắng có in dòng chữ vô tri ``I ♡ SHAKE'' (Tôi yêu sữa lắc) được đội lố bịch đè lên trên chiếc mũ lưỡi trai bóng chày quen thuộc của cô nàng. Vị ``anh hùng cứu thế'' đang đứng thẫn thờ cạnh một dây chuyền sản xuất đồ uống.

Một dòng chữ lớn bằng phông chữ công nghiệp cục súc hiện ra chình ình trên màn hình: ``Nhà máy sữa lắc số 12''.

Subaru trên màn hình đang lơ đãng dùng tay đẩy đẩy những ly nhựa đựng sữa lắc chạy trên dây chuyền tự động, vẻ mặt bơ phờ tràn ngập sự chán nản của một kẻ làm công ăn lương.

\img{4-4e}

\il

Nhìn thấy cảnh tượng đó, Ryuuhou sặc luôn bã kẹo cao su, đập tay xuống đùi cười sặc sụa: ``Vãi cả mho! Anh hùng cứu thế cái gì mà lại đi làm công nhân nhà máy vắt sữa lắc thế này!? Định mệnh vinh danh cái quần què gì nghe nghiệt ngã hãm tài thế! Thằng nào viết kịch bản chắc đang hít cỏ à!''

Kuon cũng ngã ngửa ra sau ghế sô pha, dở khóc dở cười vuốt mặt: ``Anh cũng chịu thua với độ vô tri của cái game này rồi. Giấc mộng anh hùng vỡ nát tươm con mẹ hàng lươn luôn.''

Pekora ngồi ở góc thì vẫn lơ mơ, đôi mắt cá chết lờ đờ chớp chớp lẩm bẩm: ``Peko\dots\ Mình đang thèm ăn một củ cà rốt giòn rụm, peko\dots''

\il

Trên màn hình, Subaru bực dọc quay sang cau có lườm mụ giám sát viên của nhà máy - người có tạo hình không ai khác chính là Marine - và lớn tiếng chất vấn: ``\textbf{Này bà Giám sát! Tại sao tôi, đường đường là một vị anh hùng vĩ đại được thần linh chọn lọc, lại phải bán mặt cho đất bán lưng cho trời làm cái công việc vất vả tay chân trong cái nhà máy rẻ rách này hả? Tôi không hề có thời gian rảnh rỗi để lo cho mấy cái cốc sữa lắc vớ vẩn này đâu nhé!}''

\img{4-4f}

Marine, lúc này cũng đang khoác bộ đồng phục công nhân tương tự như Subaru nhưng được thiết kế trễ nải quyến rũ hơn, vừa vô tư ngoáy ``phao câu'' nhún nhảy theo điệu nhạc xập xình ầm ĩ trong nhà máy, vừa nhếch mép cười đáp trả một cách vô cùng ngẫu hứng: ``Thì do luật pháp của nhà vua ban hành bảo vậy, nên bọn mình cứ thế mà tuân lệnh nhắm mắt làm thôi.''

\il

Ryuuhou nhăn mặt, chỉ tay thẳng vào màn hình chửi đổng: ``Đồng cây mát mẻ, bà nội Hải tặc này bị lậm vai diễn à? Nước đến chân rồi mà vẫn còn đứng đó múa may quay cuồng được!''

Kuon vò đầu thở dài: ``Đúng là cái công ty toàn những thành phần bất hảo. Vị anh hùng gánh vác thế giới lại bị 'phốt' chèn ép bởi chính nhà vua cơ đấy.''

\il

Nàng ``Anh hùng'' Subaru thở dài thườn thượt, ánh mắt tràn đầy sự thất vọng tột độ với cuộc sống: ``Cái gì mà định mệnh được lựa chọn chứ\dots\ định mệnh cái quần què ấy! Đúng là một cú lừa thế kỷ!''

\il

Kuon huých cùi chỏ sang cô nàng Thỏ xanh đang ôm tay cầm, hỏi dò: ``Này, Peko-chan, em thấy sao hả? Vị anh hùng trong mơ của em đang cắm mặt đi 'vắt sữa' kìa.''

Pekora vẫn không hề có phản ứng gì rõ rệt, đôi mắt đờ đẫn vẫn chỉ hướng về một cõi hư vô: ``Ừm\dots\ cà rốt ngon quá\dots''

Ryuuhou đảo mắt ngán ngẩm nhìn con Thỏ bên cạnh: ``Ông anh bỏ cuộc đi, con thỏ phê game này hồn bay phách lạc luôn rồi, gọi nó cũng như đàn gảy tai trâu thôi. Kiểu này thì hai anh em mình xác định phải gánh tạ gãy lưng rồi.''

Kuon gãi gãi đầu cười gượng: ``Thì em cứ ráng xem tiếp đi. Biết đâu vào sâu mấy đứa trong này lại có trò mèo gì đó tấu hài hay hay cho mình giải trí thì sao.''

Ryuuhou khoanh tay, hậm hực lầm bầm: ``Giải trí cái nỗi gì, em xem mà em bắt đầu cảm thấy tăng xông ức chế giùm cho cái con nhân vật nữ anh hùng kia rồi đấy. Game rác vãi linh hồn.''

\il

Quay trở lại trong game, Subaru nghe xong lời giải thích lười biếng của nàng Thuyền trưởng liền tỏ vẻ không hề phục tùng: ``Giám sát à, tôi không có dư dả thì giờ cho mấy thứ sữa lắc linh tinh này đâu! Trong lúc chúng ta đứng chôn chân ở đây để cãi nhau, thì cái thế lực tà ác ngoài kia đã tiếp tục bành trướng, thao túng và kiểm soát tâm trí của toàn nhân loại rồi! Chúng ta cần phải cầm vũ khí hành động ngay lập tức!''

Tên giám sát viên Marine lập tức chỉ tay thẳng vào mặt cô nàng Tomboy, đáp trả lại bằng một thái độ vô cùng nghiêm túc nhưng lại sặc mùi châm biếm: ``Cô dám mở miệng gọi sữa lắc là thứ linh tinh hả!? Cô có biết là toàn bộ nhân loại đang vô cùng thèm khát và cần đến những ly sữa lắc ngọt ngào đến thế nào trong cái thời buổi chiến tranh hỗn loạn này không!? Sữa lắc chính là cứu cánh, là niềm hy vọng duy nhất còn sót lại của loài người đấy!''

\il

Kuon tặc lưỡi: ``Nói thật là nó linh tinh thật mà. Thằng điên nào lại cần uống sữa lắc để lấy sức đi cứu thế giới chứ?''

Ryuuhou đập mạnh tay xuống ghế sofa, cười cợt nhả: ``Thề luôn, tư duy của con mụ giám sát viên này ảo ma Canada thật sự. Thế giới sắp bị hủy diệt đến nơi rồi mà bả lại đi ép người ta làm công nhân vắt sữa lắc. Cạn cmn lời với logic game Nhật! Phải tay em mà là con anh hùng kia thì em vác gươm tẩn cho con mụ giám sát một trận nhừ tử rồi chuồn luôn cho rảnh nợ!''

\il

Nàng Tomboy vẫn đang phải cố gắng gồng mình kìm nén cơn tức giận trào dâng, cãi lý lại: ``Nhưng thưa giám sát, với tình hình cấp bách hiện tại, ưu tiên hàng đầu của chúng ta là phải tập trung toàn bộ lực lượng vào việc tiêu diệt thế lực tà ác, chứ không phải\dots\ đứng đây làm sữa lắc!''

Marine lập tức nhướng mày, hùng hổ tiến tới đứng sát rạt vào mặt cô nàng và bồi thêm một câu vặn vẹo: ``Tiêu diệt thế lực tà ác sao? Được thôi, vậy thì cô hãy thử dùng cái bộ não anh hùng của mình để trả lời câu hỏi này của tôi đi: giữa số người đang bị bọn chúng bắt giữ tra tấn và số người đang chết khát cần bú sữa lắc để sống lay lắt qua ngày, cái nào mang tính sống còn quan trọng hơn hả?''

Subaru ngập ngừng, miệng há hốc không biết phải phản biện lại cái lý lẽ cùn đó thế nào, bối rối đáp: ``T-Tôi\dots\ tôi cũng không biết nữa\dots''

Chỉ chờ có thế, tên giám sát viên quỷ quyệt kia liền nhấc bổng ``vị anh hùng'' nhỏ bé lên trời, tiếp tục dùng quyền lực cấp trên chèn ép thêm bằng một câu chốt hạ phũ phàng: ``Vậy thì chốt lại, nhiệm vụ sản xuất sữa lắc để cung cấp năng lượng cho mọi người chẳng phải mới là việc hệ trọng nên được ưu tiên hàng đầu hay sao?'' Nói rồi, ả thẳng tay ném mạnh Subaru ngã nhào xuống một đống vỏ ly sữa lắc chưa kịp đóng gói dưới sàn.

\img{4-4g}

\il

Ryuuhou vỗ đùi cười sằng sặc: ``Đỉnh cao vãi chưởng! Con mụ giám sát viên này thao túng tâm lý mỏ hỗn bá đạo quá! Đỉnh cao của nghệ thuật bóc lột tư bản là đây chứ đâu!''

Kuon xoa xoa trán: ``Thật sự là não anh sắp bị nhũn ra vì không thể hiểu nổi cái mạch logic rác rưởi của cái game này nữa rồi. Trà sữa lắc với thế giới hậu tận thế diệt vong thì có cái sợi dây liên kết quái gì với nhau!?''

Cô em gái liếc nhìn chiếc máy console xịt khói một cách vô cùng thù hằn: ``Em thề luôn là em đang rất muốn tìm ra cái info thằng dev thiết kế ra cái kịch bản game củ chuối này để\dots\ tẩn cho nó một trận nên hồn đấy.''

\il

Sau cú ném điếng người, ``vị anh hùng'' Subaru lồm cồm bò dậy từ đống ly nhựa, xoa xoa cái lưng ê ẩm, miệng không ngừng than vãn trong sự cay cú tột cùng: ``Mẹ kiếp, mình thực sự muốn nộp đơn xin đổi nghề luôn quá! Tự dưng đi phí thời gian dây dưa cãi nhau với bà giám sát viên hắc ám này từ nãy tới giờ để làm cái quạc gì không biết\dots\ Trong khi nhân loại ngoài kia hiện đang lâm nguy kêu cứu!''

Dẫu trong lòng tức giận trào sôi là thế, nhưng cuối cùng dưới áp lực của đồng tiền bát gạo, cô nàng vẫn phải ngậm đắng nuốt cay quay trở lại vị trí dây chuyền sản xuất, tiếp tục làm cái công việc vắt sữa nhàm chán đến phát điên.

\ct

Cày cuốc được một lúc lâu, ánh mắt mệt mỏi của Subaru bỗng liếc thấy một bóng người vô cùng quen thuộc đang đứng làm việc ở phía đối diện của dây chuyền. Hóa ra đó lại là một cô nàng nhân viên ``bất đắc dĩ'' khác bị hút vào game, nàng Ma cà rồng Mel.

Subaru lập tức lóe lên một tia hy vọng, nhủ thầm trong đầu: ``\textit{Khoan đã, nhìn cái tạo hình thanh lịch kia\dots\ Không lẽ đó chính là vị Hiền nhân Mel trong truyền thuyết sao? Nhưng tại sao một nhân vật có tiếng tăm như chị ấy lại cũng bị bắt đi làm công nhân ở cái xưởng sữa vắt kiệt sức lao động này nữa!?}''

\img{4-4h}

\il

Ryuuhou nhún vai nhai kẹo chóp chép, châm chọc: ``Ờ thì kịch bản game rác nó nhét bừa NPC vào thế đấy. Tìm logic ở đây thì có mà đổ thóc giống ra mà ăn.''

\il

Nàng Tomboy thấy đây chính là cơ hội vàng ngàn năm có một để thoát khỏi cái công việc tù đày này, liền nảy ra một ý tưởng đào tẩu. Cô hít một hơi thật sâu, hai bàn tay mạnh mẽ nắm chặt lại như thể đã hạ quyết tâm lớn: ``Được rồi! Nếu có Hiền nhân giúp sức, hai chúng ta chắc chắn có thể liên thủ hợp tác để trốn thoát khỏi cái nhà máy quỷ tha ma bắt này được!''

Cô gọi với sang nàng Dơi đang chăm chỉ làm việc một cách vô cùng vui vẻ và tận tụy ở đằng kia: ``\textbf{Bậc Hiền nhân đáng kính ơi!!}''

Nghe tiếng gọi réo rắt, Mel từ từ quay đầu lại. Với một nụ cười hiền hậu, dịu dàng tỏa nắng thường thấy, nàng Ma cà rồng cất giọng hỏi han: ``Ồ, ta nhìn nhầm hay sao, chả phải đó chính là vị anh hùng vĩ đại Subaru đấy sao?''

Như người chết đuối vớ được cọc, Subaru liền gật đầu lia lịa chào hỏi: ``Dạ đúng là con đây! Lâu quá rồi không gặp ngài, thưa Hiền nhân!'' Màn chào hỏi nhiệt tình đến mức làm rớt luôn cả cái mũ công nhân màu trắng xuống sàn.

``Đúng ha, tính ra cũng phải một khoảng thời gian dài đằng đẵng rồi tụi mình chưa gặp lại nhau. Dạo này công việc cày cuốc của con ở cái nhà máy này thế nào rồi?'' Mel ân cần hỏi.

Subaru nghe hỏi thăm liền tủi thân méo xệch cả mặt, bắt đầu thao thao bất tuyệt kể lể nỗi khổ tâm với cô nàng: ``Dạ khổ lắm ngài ơi, con vừa mới bị cái tên giám sát viên độc ác kia mắng chửi té tát đổ lên đầu xong\dots\ Cô ta thậm chí còn bạo lực đến mức dùng võ thuật bẻ gãy gập cả cái lưng tội nghiệp của con luôn đây này!''

Nghe những lời tâm sự đầy cay đắng của ``vị anh hùng'', Mel liền đưa tay đặt lên cằm, gật gù ra vẻ thông cảm như một nhà hiền triết vô cùng thấu tình đạt lý. Sau một hồi nhắm mắt suy tư sâu sắc, nàng nghiêm túc phán một câu xanh rờn: ``Ta hoàn toàn thấu hiểu nỗi khổ của con. Thế nhưng mà con ạ, con cũng đừng nên mang ác cảm nghĩ xấu về cô giám sát viên nghiêm khắc kia, vì sâu thẳm trong thâm tâm, ta tin chắc là cô ấy làm vậy cũng chỉ là vì muốn tốt cho sự nghiệp của con thôi mà.''

\img{4-4i}

\il

Kuon nghe xong câu thoại liền phụt cười sặc sụa: ``Ôi Mel ơi là Mel, anh thấy em nhập vai hiền hòa thánh thiện quá lố rồi đó\dots\ Nói thật mất lòng nha, ngoài đời đào đâu ra cái thể loại sếp bạo hành nhân viên vì muốn tốt cho nhân viên cơ chứ.''

Ryuuhou gật đầu lia lịa tán thành: ``Chuẩn đét! Thằng sếp giám sát viên nào mỏ hỗn kêu `muốn tốt cho nhân viên' thì một trăm phần trăm là mồm điêu xạo chó bóc lột sức lao động. Cú lừa kinh điển của bọn tư bản! Hơn nữa, nhìn cái mặt con mụ Marine kia mà bảo là người tốt thì em thà tin mẹ nó vào việc thằng trùm cuối phản diện là người hay đi phát gạo từ thiện còn có lý hơn!''

\il

Thế nhưng, Subaru trên màn hình lại hoàn toàn không có khả năng nhận thức được sự ngây thơ quá đáng trong lời khuyên đó. Cô nàng chỉ biết nhìn vị Hiền nhân Mel với một ánh mắt bừng sáng rực rỡ, dường như những lời động viên sáo rỗng ấy đã tiếp thêm cho cô một luồng sinh khí, một niềm hy vọng mới mẻ để tiếp tục chiến đấu\dots\ à không, tiếp tục vắt sữa lắc.

Subaru khẽ thở dài nhẹ nhõm, rồi nhanh chóng nảy ra một ý tưởng rủ rê để bày tỏ lòng biết ơn sự chỉ bảo của bậc tiền bối. Cô nghiêng đầu, cố gắng nặn ra một nụ cười tươi tắn nhất có thể: ``À, đúng rồi Hiền nhân ơi! Tối nay sau khi tan ca làm mệt nhọc, hai chúng ta cùng nhau rủ đi ăn tối giải khuây nhé? Chầu này cứ để con bao trọn gói cho!''

Đôi mắt màu hổ phách của nàng Dơi khẽ ánh lên một tia ngạc nhiên thích thú, nhưng rồi ngay lập tức đôi môi lại cong lên thành một nụ cười vô cùng dịu dàng, trìu mến. Nhận được sự đồng ý, Subaru bỗng chốc trở nên cực kỳ phấn khích, hai tay nắm chặt lại đưa lên cao, reo hò ầm ĩ hệt như một đứa trẻ được nhận quà: ``Hể? Ngài đồng ý đi thật ạ!? Hoan hô! Vậy thì quyết định vậy nhé, sau giờ tan ca là ta triển luôn\~{}!''

\il

Kuon cười bò: ``Anh nghĩ cái thằng cha thiết kế game này nên đổi hẳn tên game thành trò `giả lập làm culi công ty tư bản' thì sát với thực tế hơn là cái mác RPG cứu thế\dots''

Ryuuhou vung tay ngán ngẩm tột độ: ``Thôi thôi ông anh đừng có nói nữa. Game rác thì muôn đời là game rác, nó quên mịa mất luôn cái cốt truyện chính là đi giết quái cứu thế giới rồi. Toàn đưa mấy cái kịch bản xàm lờ câu giờ vớ vẩn.''

\il

Mel nở nụ cười tươi rói, cất giọng đầy vẻ khích lệ động viên tinh thần: ``Chỉ còn khoảng bốn tiếng đồng hồ tăng ca nữa là tụi mình được nghỉ rồi. Hai ta cùng nhau cắn răng cố gắng cày nốt cho xong nhé?'' 

Subaru nghe thấy vị Hiền nhân kính yêu nói vậy, dẫu cho trong lòng vẫn còn chất chứa chút bực bội ấm ức với công việc nhàm chán vắt kiệt sức hiện tại, nhưng cũng chẳng nỡ mở lời từ chối trước thái độ nhiệt tình và thánh thiện quá mức của ``Bậc Hiền nhân''. Cô nàng dứt khoát gật đầu thật mạnh, đáp lời đầy quyết tâm rực lửa: ``Dạ vâng ạ thưa ngài! Con xin cảm ơn ngài rất nhiều vì đã khai sáng!''

\il

Ryuuhou lắc đầu xua tay bỉ bôi: ``Nói thật nhé, em mà vô tình bị kẹt vào cái hoàn cảnh oái oăm đó là em AFK (treo máy) thoát  game luôn từ vòng gửi xe rồi đấy. Chơi cái thể loại game trầm cảm này không khéo hệ thần kinh còn bị tàn phá nặng nề hơn cả ngoài đời đi làm thực tế.''

Kuon cười lớn đồng cảm: ``Haha, thì cứ coi như đây là một cái minigame thử thách lòng kiên nhẫn của người chơi đi em. Nhưng thú thật là anh cũng bắt đầu thấy tò mò không biết cái đoạn `giải cứu thế giới' cốt lõi cuối cùng nó sẽ lố lăng và xàm xí đến cái mức độ nào nữa đây.''

\il

Trong lúc đang hừng hực khí thế làm việc, bất chợt, Subaru nhận ra có một điều gì đó vô cùng bất thường. Cô nàng liếc mắt nhìn lên chiếc đồng hồ treo tường lớn của công ty, rồi ngay lập tức giật mình hét toáng lên đầy bức xúc: ``\textbf{Ủa khoan đã? Rõ ràng là mình nhớ ca làm việc chính thức của mình đã kết thúc từ nãy rồi cơ mà! Giờ này là giờ về rồi! Thế thì thôi dẹp đi, cắp túi xách đi về lẹ nào!}''

Cô nàng hớn hở vứt bỏ đồ đạc, quay lưng chuẩn bị sẵn sàng tư thế để chuồn lẹ khỏi cái nhà máy địa ngục. Nhưng ngay khi gót chân vừa mới cất bước, một giọng nói cực kỳ hắc ám, đầy uy quyền bỗng vang lên từ phía sau gáy như một gáo nước lạnh tát thẳng vào mặt: ``Ấy ấy, từ từ đã cô em, bình tĩnh nào. Cô vẫn còn cả một núi công việc đang chất đống chờ xử lý kia kìa!''

\img{4-4j}

Subaru bực dọc quay phắt người lại, đập vào mắt là gương mặt hắc ám của Marine. Lúc này, mụ giám sát viên còn đang cố tình đeo thêm một chiếc kính râm màu đen ngòm - tạo hình nhìn ngứa mắt y hệt như mấy tay trùm phản diện buôn lậu trong mấy bộ phim hành động kinh phí thấp rẻ tiền.

``Nữ anh hùng'' lập tức cau mày nhăn nhó, hùng hổ bật lại ngay tắp lự: ``Cái gì cơ chứ!? Cô đùa tôi đấy à! Luật lao động của nhà nước Nhật Bản đã quy định rành rành thời gian làm việc tối đa của công nhân chỉ có đúng tám giờ một ngày thôi nhé! Tôi đã làm quá giờ rồi, tôi hoàn toàn có quyền hợp pháp được nghỉ ngơi!''

\il

Ryuuhou nhếch mép cười khẩy bóc mẽ: ``Nói thế thôi chứ bên Nhật cái văn hóa ép nhân viên tăng ca OT nó phát triển ghê gớm lắm. Kiểu cái việc tăng ca sấp mặt đến đêm muộn là luật ngầm mẹ nó rồi, méo ai thoát được đâu.''

Kuon gật gù đồng tình: ``Thế thì đẻ ra ba cái thứ điều luật giới hạn thời gian làm việc đó để làm cái kiểng gì cho phí giấy mực nhỉ?''

\il

Bỏ ngoài tai mọi lý lẽ, nàng Vịt vẫn tiếp tục gân cổ lên tranh luận nảy lửa đòi quyền lợi: ``Luật vẫn mãi là luật nhé! Tôi cóc cần quan tâm cô ngụy biện nói cái gì! Tôi có quyền con người là được phép nghỉ ngơi!''

Marine, vẫn với dáng vẻ vô cùng bình thản đến đáng ghét, buông ra những lời lẽ sắc bén như dao lam đâm thẳng vào tim, chỉ khẽ nhún vai và đáp lại lạnh tanh: ``Ủa cô em bị mất trí nhớ rồi sao? Ai nói với cô cái vùng đất hoang tàn này là nước Nhật cơ chứ? Ở đây làm gì có cái luật đó.''

Câu nói phũ phàng ấy khiến Subaru lập tức đứng hình như tượng đá, đôi mắt mở to hết cỡ ngạc nhiên tột độ. Nhưng cái tên giám sát viên tư bản kia hoàn toàn không để cho đối phương có cơ hội kịp định thần phản ứng, ả tiếp tục bồi thêm một đòn tấn công tâm lý bằng một ``sự thật'' còn đáng sợ và rùng rợn hơn gấp vạn lần: ``Với cả, nói cho cô em biết để mở mang tầm mắt nhé, ở cái nhà máy này, thời gian làm việc tối đa cho phép là\dots\ 700 tiếng cơ. Chúc cô em có những giây phút tăng ca thật vui vẻ nhé\~{}''

Nói xong câu nói sặc mùi máu lạnh, cô nàng giám sát viên thản nhiên quay lưng bỏ đi một nước, để mặc lại nàng Vịt đáng thương đứng chết trân tại chỗ, há hốc miệng không nói nên lời.

Cả Subaru trong màn hình, lẫn Kuon và Ryuuhou đang cầm tay cầm ở bên ngoài thế giới thực, đều không kìm chế nổi sự phẫn nộ, cùng đồng thanh hét lên phẫn uất: \textbf{Thế này thì khác éo gì là bóc lột sức lao động của người ta đến tận xương tận tủy cơ chứ!!}

\il

Cô em gái nghiến răng ken két, chửi thề: ``700 tiếng tăng ca!? ĐM thằng dev rác rưởi nào làm ra cái game vô nhân tính đốn mạt này thế. Nó nghĩ ai cũng là trâu ngựa chắc?''

Cậu Tổng thấy em gái mình nổi cáu, liền vỗ vỗ vai em gái xoa dịu: ``Thôi bớt nóng đi cô nương. Anh cá là nó cố tình viết ra cái luật xàm lờ này cốt chỉ để gây ức chế tâm lý cho người chơi thôi. Mấy cái thể loại game nhảm nhí này thường thì kịch bản càng vô lý, càng ức chế thì lại càng hút khách mua về chửi mà.''

\il

Cuối cùng, khi giọt nước đã tràn ly, không thể nào tiếp tục chịu đựng nổi sự áp bức bóc lột dã man của mụ giám sát Marine được nữa, Subaru giận dữ đến tột đỉnh, hét lớn vang vọng cả nhà máy: ``\textbf{Bà mày méo thèm làm cái công việc nô lệ này nữa!}''

Nàng ``Hải tặc'' nghe tiếng chửi phản kháng ngay lập tức quay phắt người lại. Ả từ từ gỡ bỏ chiếc kính râm xuống, lăm lăm nheo đôi mắt nhìn chằm chằm vào ``vị anh hùng'' với vẻ mặt đầy rẫy sự nguy hiểm sát khí và gầm gừ hỏi: ``Cô vừa sủa cái gì cơ? Có giỏi thì mở cái mõm ra lặp lại câu vừa nói xem nào!?''

Chưa kịp để cho Subaru phản ứng hay thủ thế phòng vệ, nàng Vịt đã bị mụ giám sát lao tới tóm chặt lấy, bạo lực bẻ quặt người cô gập hẳn xuống một cách vô cùng tàn nhẫn không chút thương tiếc. Cô giám sát viên hung hăng gằn giọng đe dọa vào lỗ tai nạn nhân: ``Cô nghĩ rằng một kẻ rác rưởi như cô có thể tồn tại và làm việc cho cái công ty vĩ đại này với cái suy nghĩ chống đối phản loạn đó hả? Tốt nhất là ngậm cái mõm lại và tiếp tục cắm đầu vào làm việc cật lực ở cái xưởng này thêm 900 tiếng đồng hồ nữa cho tôi! Cấm cãi!''

\img{4-4k}

\il

Ryuuhou trợn mắt, chỉ tay vào màn hình chửi ầm lên: ``Cái qué gì thế!? Ép người quá đáng vãi chưởng thế này\dots\ Mà cái con anh hùng này phế vật yếu đuối thật sự! Vùng lên tẩn cho con mụ giám sát một trận đi chứ con Vịt kia!''

Kuon vò đầu bứt tóc, bóp trán than vãn: ``Thề luôn là sức chịu đựng của anh đang cạn kiệt và anh bắt đầu thấy cực kỳ ghét cái game này rồi đấy. Nó đang cố tình bào mòn và thử thách giới hạn sự kiên nhẫn của người chơi.''

Cô em gái hừ lạnh: ``Anh tưởng mình anh cay chắc? Cái này mà không phải vì nể tình ông anh năn nỉ rủ chơi để giải cứu đám nhân viên ngốc nghếch kia thoát ra ngoài thì bà mày đã vác búa đập nát cái tivi này từ lâu rồi.''

\il

Subaru hét lên những tiếng thảm thiết đầy đau đớn khi bị Marine dùng lực bẻ quặt cột sống lưng: ``\textbf{Aaa! Đau! Đau quá gãy lưng tôi mất! Được rồi, tôi biết sai rồi! Xin giám sát hãy dừng tay lại đi mà!}''

Sau khi đã hoàn thành màn trừng phạt bạo lực răn đe kẻ chống đối, Marine lạnh lùng hất đối phương ra, tự bẻ khớp vai mình rắc rắc như thể vừa mới chiến thắng một trận quyết đấu sinh tử. Trước khi kiêu hãnh rời gót khỏi xưởng sản xuất, cô ả còn không quên buông lại một câu răn đe đầy máu lạnh: ``Khôn hồn biết điều nghe lời thế là rất tốt. Lo mà cắm đầu vào làm việc cho cẩn thận và đàng hoàng vào đấy. Chào nhé lũ nô lệ!'' rồi nhẫn tâm sải bước rời đi.

Subaru nằm gục rã rời xõa xượi trên nền đất lạnh lẽo, cơ thể đau nhức, tức tối và hoàn toàn kiệt sức. Cô nàng nghiến răng trèo trẹo, dùng tay đấm mạnh xuống mặt sàn lẩm bẩm đầy thù hằn: ``Được thôi, cứ chờ đấy con mụ già ác độc, nếu cô đã ép tôi đến bước đường cùng thế này thì tôi thề sẽ cày cuốc để soán luôn cái chức giám sát viên của cô! Để xem rốt cuộc ai mới là kẻ cười cuối cùng!''

\img{4-4l}

Cô nàng lồm cồm vùng đứng dậy, ánh mắt lóe lên một tia quyết tâm báo thù mãnh liệt. Nhưng trớ trêu thay, chính vào cái khoảnh khắc sục sôi ý chí leo rank công sở đó, cái giấc mơ vĩ đại về sứ mệnh trở thành một vị anh hùng giải cứu thế giới của cô nàng đã chính thức bị ném thẳng vào dĩ vãng lãng quên.

\ct

Trong khi đó, ở một không gian chiều không gian đen tối hắc ám khác, tên cầm đầu của ``thế lực tà ác'' Choco đang đứng khoanh hai tay trước ngực ngạo nghễ ngay giữa trung tâm một căn phòng u ám đang bốc cháy hừng hực những ngọn lửa đỏ rực. Trái với phong thái của một trùm cuối nguy hiểm, vẻ mặt của ả lại đang lộ rõ sự chán nản, ngáp ngắn ngáp dài chờ đợi.

``Rốt cuộc là bọn anh hùng rơm đó đang làm cái trò trống gì mà lề mề đi lâu đến thế không biết? Chẳng lẽ đoạn đường tìm đến cái hang ổ này lại giăng bẫy khó khăn quá sức với lũ hạ đẳng đó lắm sao?''

\img{4-4m}

\il

Ryuuhou bật cười lớn sảng khoái vỗ đùi: ``Em lạy chúa, em thực sự xin giơ hai tay đầu hàng với cái kịch bản này luôn rồi. Cày game bao nhiêu năm nay, em thề là chưa từng thấy cái tựa game RPG nào mà nhân vật anh hùng chính lại u mê làm culi đến mức quên cmn luôn cả nhiệm vụ đi đánh boss cứu thế giới như thế này.''

Kuon thở dài não nề, ngón tay bóp chặt lấy cái tay cầm, giọng trầm hẳn xuống đượm mùi chán nản: ``Thôi, anh xin em bớt phán xét đi, tập trung try-hard phá đảo nhanh gọn lẹ cái con game củ chuối này rồi còn đi về. Anh không muốn lãng phí thêm một giây một phút cuộc đời nào cho cái đống rác rưởi này nữa.''

Ryuuhou gật đầu đánh rụp: ``Chuẩn. Tốc chiến tốc thắng kết thúc cái mớ hỗn độn này thôi.''

Kuon quay sang nhìn con Thỏ đang ngồi thất thần: ``À, và Peko-chan ơi? Em làm ơn cũng tập trung điều khiển nhân vật phụ anh chị một tay nhé? Xin em đừng có đi loăng quăng bóp team báo hại cả nhóm cày lại từ đầu đấy.''

Pekora chớp chớp mắt cá chết: ``P-Peko\dots\ Cà rốt\dots\ Peko\dots'' rồi cũng gật gù trong vô thức.

Ryuuhou nhìn con thỏ ngán ngẩm: ``\dots Thôi thì cứ tạm cho là nó đã lọt tai được lời ông anh nói đi vậy. Vào việc!''

\il

Cả ba người bắt đầu chấn chỉnh tư thế tiếp tục dán mắt vào màn hình cày game, nhưng cái thứ cảm giác háo hức, hồi hộp khám phá ban đầu giờ đây đã hoàn toàn tan biến sạch sẽ không còn một dấu vết. Thay thế vào đó là một sự cam chịu, cày cuốc trong sự cáu bẳn một cách vô cùng kỳ lạ, cứ như thể chính bản thân những người đang ngồi cầm tay cầm ở ngoài đời thực cũng đang bị mắc kẹt tâm trí vào cái guồng quay lao động của xưởng sữa lắc đầy phi lý, điên rồ kia vậy.

%==========================================TRUYỆN PHỤ=========================================
\omk

Toàn bộ những phân cảnh khôi hài mà Kuon và Ryuuhou vừa phải chứng kiến từ góc nhìn của bộ ba nhân vật Subaru, Mel và Marine ở phần trên thực chất mới chỉ là phần mở đầu (prologue) làm quen của cái trò chơi rác rưởi đó. Về sau, khi cốt truyện dần hé mở, họ mới tá hỏa ngã ngửa ra khi biết được một cú ``quay xe'' kịch bản kinh hoàng không ai có thể lường trước được: hóa ra, chủ nhân quyền lực đứng đằng sau giật dây điều hành toàn bộ công ty sản xuất sữa lắc bóc lột đó lại chính là Choco, kẻ cũng đồng thời là trùm cuối đứng đầu thế lực hắc ám tà ác đã nhẫn tâm xâm chiếm, tàn phá gần như toàn bộ nền văn minh nhân loại.

Tất cả mọi mớ hỗn độn đang diễn ra, từ việc thao túng xây dựng công ty sữa lắc đày đọa công nhân, cho đến những bộ quy định lao động khắc nghiệt, bóc lột sức người tới mức vô lý được lập trình trong trò chơi, hóa ra đều là một phần nằm trong kế hoạch hiểm độc đã được tính toán kỹ lưỡng của ả ta (hoặc ít nhất, là của nhân vật phản diện xảo quyệt mà nàng Y tá xui xẻo đang phải đóng vai).

``Cái thể loại kịch bản méo gì thế này!? Kẻ đứng ra điều hành cả cái tập đoàn bóc lột sức lao động địa ngục đó lại là\dots\ Choco-san sao!? Trùng hợp đến mức này thì đúng là chịu thua!'' cậu Tổng quản lý há hốc mồm thốt lên đầy ngỡ ngàng khi tình tiết được phơi bày.

Nhưng đối với một game thủ cày game nhập vai lão làng đầy mình kinh nghiệm thực chiến như Ryuuhou, cô bé lại tỏ ra không mấy ngạc nhiên hay bị sốc trước cú ``quay xe'' rẻ tiền này: ``Cũng dễ đoán thôi. Thảo nào mà nhân vật anh hùng của tụi mình từ đầu game tới giờ bị hành cho tơi bời hoa lá không ngóc đầu lên nổi\dots\ Hóa ra kẻ đứng sau giật dây bóc lột lại là ả boss cuối! Kịch bản đúng kiểu rác của rác!''

Trên màn hình ti-vi, nhân vật nữ phản diện Choco cuối cùng cũng chịu xuất hiện lộ diện. Ả ta bước ra với dáng điệu lả lơi, uốn éo đầy kiêu ngạo, trên tay cầm hờ hững một ly rượu vang đỏ như máu lắc nhẹ: ``Hoan nghênh các ngươi, những con thiêu thân ngu ngốc đã lết xác tìm được đường đến tận đây. Nhưng đáng tiếc thay, mọi nỗ lực cày cuốc thảm hại của các ngươi từ trước tới giờ chung quy đều trở nên vô ích mà thôi.''

Subaru, nhân vật lúc này đang được Ryuuhou điều khiển vô cùng cục súc, thốt lên phẫn nộ: ``Mụ điên kia, cô đang nói cái gì thế hả!?''

Choco khẽ ngửa đầu cười nhếch mép, cất lên tông giọng lanh lảnh, the thé đặc trưng của một ả ma nữ ranh mãnh: ``Sự thật rất đơn giản thôi. Tất cả lũ vô dụng các ngươi, từ đầu đến cuối đều chỉ là những con tốt thí mạng rẻ rách nhảy múa trên bàn cờ trò chơi do chính tay ta thiết kế ra. Từ việc thao túng xây dựng cái công ty sữa lắc cho đến việc đẻ ra mấy cái luật lệ làm việc khắc nghiệt vắt kiệt sức, tất cả đều là một phần trong kế hoạch hoàn hảo tuyệt mỹ của ta nhằm mục đích khiến cho ý chí của các ngươi bị mài mòn, thể xác kiệt quệ và cuối cùng là tự nguyện quỳ gối phục tùng ta!''

Cô ả từ từ bước những bước đi khiêu khích tiến đến sát chỗ ba nữ anh hùng đang đứng, nhìn sâu vào đôi mắt của họ mà không hề để lộ ra chút nào sự sợ hãi: ``Bây giờ, nếu biết thân biết phận thì hãy ngoan ngoãn quỳ phục rạp xuống dưới chân ta đi lũ thất bại.''

Mel, nhân vật hiền triết đang được Kuon điều khiển, vẫn cố chấp giữ thái độ ngây thơ, lên tiếng phân bua giảng đạo lý với nữ ác nhân: ``T-Tại sao cô lại có thể nhẫn tâm đối xử độc ác với nhân gian như vậy chứ?''

Trong khi đó, bản thân con bé Ryuuhou ngồi cầm tay cầm ở ngoài đời thực cũng đã bắt đầu tỏ vẻ bực bội đến cực điểm, cáu gắt đến mức dường như màng ranh giới giữa đời thực và thế giới ảo đã hoàn toàn bị xóa nhòa: ``Kêu bà mày quỳ xuống á? Con mụ điên này, tao nghĩ là cái thứ mày sắp nhận được là nguyên một cái gậy bóng chày vụt thẳng vào mồm thì đúng hơn đấy! Đứng đó mà gáy!''

Thấy con em gái bắt đầu có dấu hiệu tăng xông máu dồn lên não, Kuon vội vã đưa tay sang vỗ vỗ vai xoa dịu: ``Thôi nào, bớt nóng đi Ryuuhou ơi, hít thở sâu hạ hỏa đi\dots\ Bọn mình phải cố gắng giữ cái đầu lạnh, nghĩ cách đánh boss giải quyết êm xuôi chuyện này một cách cẩn thận\dots\ Manh động không giải quyết được gì đâu, lỡ chết lại mất công cày lại từ checkpoint thì khổ.''

Về bản chất thâm sâu của nó, toàn bộ cái tựa trò chơi quái quỷ này được hệ thống ma thuật thiết kế mô phỏng như một dạng trải nghiệm nhập vai, dựa trên chính những ước muốn, khao khát thầm kín nhất của ba cô gái Subaru, Mel, và Marine, đúng y hệt như những tâm tư nguyện vọng mà họ đã từng bộc bạch thổ lộ với Kuon trước đó.

Thế nhưng sự đời vốn lắm éo le, thay vì để họ được hóa thân trở thành những anh hùng vinh quang dệt nên huyền thoại giải cứu thế giới, hệ thống game lại biến tấu vặn vẹo ước mơ đó, đày đọa họ trở thành những công nhân nạn nhân khốn khổ bị mắc kẹt không lối thoát trong guồng quay bóc lột tàn nhẫn của chủ nghĩa tư bản. Ít nhất thì đó cũng là sự thật phũ phàng theo đúng như những gì mà lời ả Choco vừa mới tuyên bố.

Sau khi nghe xong toàn bộ những lời thú tội trơ tráo từ ả Y tá, từ cái kế hoạch thâm độc nhằm thôn tính nhân loại cho đến cả cái mưu đồ giăng bẫy guồng quay lao động bóc lột, tất cả ba nữ anh hùng đứng đó đều lộ rõ vẻ bàng hoàng, sửng sốt đến tột độ.

Ngay cả đến nhân vật Marine, người bị hệ thống cưỡng ép phân công vào vai mụ nữ giám sát viên bóc lột một cách vô cùng bất đắc dĩ, đồng thời cũng chính là nhân vật đang do cô nàng Thỏ Pekora (dù vẫn còn trong tình trạng lờ đờ) nắm quyền điều khiển, cũng phải hét lên phẫn nộ với hành động phản bội, đem con bỏ chợ của ả sếp: ``Cái gì cơ!? Chẳng lẽ tất cả những việc làm ác ôn vắt kiệt sức lao động của tôi từ trước tới nay\dots\ cũng chỉ là để phục vụ mục đích biến tôi trở thành một con tốt thí mạng cho cô lợi dụng thôi sao?''

Subaru run rẩy, đôi mắt rưng rưng nước mắt vì uất ức: ``K-Không thể nào có chuyện vô lý như thế được! Tại sao số phận lại trớ trêu và nghiệt ngã đến mức như thế này cơ chứ!? Rõ ràng là tôi chỉ muốn được trở thành một người anh hùng cứu nhân độ thế cơ mà!''

Bậc hiền nhân Mel thì chỉ biết ngập ngừng, lùi lại phía sau cất lên một giọng nói yếu ớt, run rẩy: ``Còn tôi\dots\ tôi từ đầu đến cuối cũng chỉ muốn mang lại cuộc sống hạnh phúc và bình yên cho tất cả mọi người thôi mà\dots'' Nàng ma cà rồng khẽ sụt sịt, không thể kìm nén được những giọt nước mắt lăn dài trước một thực tại quá đỗi đen tối và phũ phàng này.

``Nữ anh hùng'' Subaru không thể nào chấp nhận nổi việc niềm tin và lý tưởng của mình bị chà đạp, liền gào thét lớn vào mặt ả nữ Y tá: ``Cô làm thế này là quá quắt lắm rồi đấy! Rốt cuộc mục đích của ngươi là muốn cái gì thì cứ mở cái mõm thúi ra mà nói huỵch toẹt thẳng ra đi!''

Choco ngửa cổ cười lớn đắc thắng, cất lên một tràng giọng điệu the thé, đỏng đảnh đầy khinh miệt: ``Mục đích của ta ư? Đơn giản thôi! Ta muốn nhìn thấy tất cả lũ các ngươi từ từ chìm đắm chới với trong sự tuyệt vọng cùng cực, để rồi ngoan ngoãn cúi đầu phục vụ làm culi cho ta mãi mãi không có ngày siêu thoát!''

Từ phía bên ngoài đời thực, chứng kiến thái độ kiêu ngạo lố lăng của ả trùm cuối, cả Kuon lẫn Ryuuhou đều cùng đồng loạt đanh mặt lại, toát ra một vẻ quyết tâm cao độ hừng hực sát khí. Riêng con thỏ Pekora dù vẫn còn dặt dẹo lờ đờ vì di chứng của căn bệnh nghiện game, nhưng qua nét mặt dường như cô nàng cũng lờ mờ hiểu được mức độ căng thẳng của tình thế khốc liệt hiện tại sắp sửa nổ ra.

\il

Ryuuhou bẻ khớp tay rắc rắc, mặt hầm hầm sát khí: ``Vào việc đi ông anh! Bà mày ngứa tay lắm rồi, phải tẩn chết cụ cái con phản diện rác rưởi này cho bõ tức!''

Kuon gật đầu dứt khoát, bấm chặt tay cầm: ``Đồng ý. Lên đồ đi. Tụi mình bắt buộc phải phá đảo con game này để còn lôi cổ cả ba người bọn họ thoát về thực tại nữa!''

Pekora lẩm bẩm trong miệng, tay vẫn bấm nút theo bản năng: ``Peko\dots\ đánh boss\dots\ cứu cà rốt\dots\ Peko\dots''

\il

\ct

Phải cày cuốc sấp mặt mồ hôi cho đến tận chiều muộn ngày hôm đó, sau khi trải qua một trận chiến vô cùng dài hơi, căng thẳng và trầy trật đổ mồ hôi hột, cả nhóm rốt cuộc cũng đã đánh bại được Choco trong trận đấu trùm quyết định. 

Thú thật, trên cái hành trình đánh boss gian nan đó, đã không ít lần Kuon vì thiếu kỹ năng chơi game hoặc con thỏ phê game Pekora vì bấm lộn nút mà suýt chút nữa đã đẩy cả đội vào cửa tử báo hại công cốc, và cũng không ít lần cả ba người bị hệ thống ra đòn chí mạng của nàng Y tá kia chơi cho những vố cực đau mất gần hết sạch cây máu. Nhưng bằng một nghị lực phi thường và khả năng gánh tạ tuyệt đỉnh, Ryuuhou, một game thủ có số má vươn tầm cao thủ, rốt cuộc cũng đã kéo lê được cả hai ông bà tạ kia phá đảo thành công cái tựa trò chơi quái quỷ, ức chế nọ.

Ngay khi đoạn cắt cảnh hình ảnh kết thúc game từ từ hiện lên trên màn hình, chiếc máy chơi game console tàn tạ bỗng dưng rung lắc dữ dội rồi phát sáng rực rỡ chói lòa. Luồng ánh sáng chói mắt ấy há miệng nhả văng ba người con gái ra ngoài. Subaru, Mel và Marine cuối cùng cũng đã thoát khỏi sự giam cầm của thế giới trò chơi ảo, ngã uỵch xuống sàn nhà quay trở về với thế giới thực tại bình yên.

\img{4-4n}

Khi bộ ba nạn nhân vừa mới lồm cồm bò dậy xuất hiện trở lại, Kuon vẫn cố gắng kìm nén sự mệt mỏi, giữ vẻ mặt bình thản, mặc dù trong thâm tâm cậu vẫn còn cay cú ghim thù sâu sắc sau khi vừa phải cày nát cái trò chơi rác rưởi do quỷ dữ lập trình đó: ``Mừng mọi người đã sống sót quay trở lại bình an.''

Ryuuhou thì đứng dậy vươn vai kêu răng rắc, toàn thân con bé mệt nhoài, rã rời sau hàng giờ đồng hồ ngồi dán mắt ``dùi mài kinh sử'' try-hard trước màn hình vô tuyến. Dù vậy, nó vẫn cố nhếch mép nở một nụ cười mãn nguyện của kẻ chiến thắng: ``Phù, gánh gãy cả lưng! Tóm lại là vụ cứu người coi như đã giải quyết xong xuôi êm đẹp rồi. Giờ thì chỉ còn sót lại đúng một việc cuối cùng cần phải làm thôi\dots'' Ryuuhou quay phắt sang trừng mắt nhìn Kuon và hỏi với chất giọng sát thủ lạnh tanh: ``\dots Việc bây giờ là bpkm mình phải đi truy lùng ngay cái info của cái thằng dev rác rưởi, trẻ trâu nào đã dám lập trình ra cái con game hãm lờ ức chế kia để tẩn cho nó một trận thì mới trôi cục tức này được.''

Cậu Tổng quản lý chỉ thở dài không đáp, liền im lặng đưa ngón tay chỉ thẳng vào chiếc máy chơi game console trên kệ, lúc này đã bị bốc khói đen thui và cháy khét lẹt ở phần viền nhựa, rồi cất giọng nhẹ tênh đáp lại: ``Thì thủ phạm làm ra cái con game đó chính là nó đấy em.''

Không cần chờ đợi thêm bất kỳ một giây phút giải thích nào nữa, Ryuuhou lập tức lao nhanh ra mở tủ đồ nghề của phòng nghỉ, lôi ra một cây gậy bóng chày bằng nhôm cứng cáp (mà thực ra chính Kuon cũng chẳng hiểu tại sao ở trong văn phòng công ty thần tượng lại có tàng trữ một cái vũ khí sát thương cao như vậy). Con bé nghiến răng lao đến với vận tốc bàn thờ, vung gậy đập giáng thẳng một cú chí mạng vào chiếc máy chơi game. Lực đập mạnh đến nỗi lớp vỏ nhựa cứng vỡ toác làm đôi nát bét, linh kiện văng tung tóe khắp sàn.

Nhìn con em gái đang hóa điên như một con thú dữ, điên cuồng vung gậy phá hủy cái máy chơi điện tử để xả giận, Kuon chỉ biết lùi lại phía sau vài bước để tránh miểng văng trúng, đứng nhìn cảnh tượng bạo lực không hề chớp mắt: ``Anh nghĩ là phen này anh mày kiểu gì cũng phải móc tiền túi ra đền bù thiệt hại cơ sở vật chất cho công ty mất thôi.''

``Tiền bạc thì đi cày OT là có thể kiếm lại được! Nhưng cái mớ thời gian thanh xuân quý báu cùng với mấy tiếng đồng hồ tao phải chịu đựng cay cú, rước bực vào người khi bị hành hạ trong cái game rác rưởi này thì\dots'' Ryuuhou vừa chửi đổng vừa vung gậy đập bồi thêm vài nhát sấm sét nữa vào tàn tích chiếc máy, lầm bầm chửi qua từng nhịp thở dốc hổn hển: ``Éo bao giờ lấy lại được. \textbf{Éo bao giờ luôn nghe chưa cái đồ sắt vụn rác rưởi này!!}''

Cả ba cô gái Subaru, Mel và Marine vẫn còn đang trong trạng thái ngơ ngác, chưa kịp hoàn hồn lại sau cú sốc khi vừa mới thoát ra khỏi thế giới trò chơi ảo. Vừa thấy cảnh tượng Ryuuhou hung hãn đập nát bét chiếc máy thành một đống sắt vụn không thương tiếc, cả ba liền kinh hãi lùi dạt hết ra góc tường.

Tuy nhiên, Subaru bỗng dưng lại giơ nắm đấm lên trời, hòa nhịp cùng cơn tức giận hò reo cổ vũ nhiệt liệt cho cô bé mắt hai màu nọ: ``Đập mạnh tay nữa lên đi Ryuuhou-chan! Đập nát bét nó ra cho bõ tức!''

Bậc hiền nhân Mel thì vẫn giữ bản tính nhút nhát, rụt rè tiến lên khuyên nhủ: ``E-Em hoàn toàn có thể thấu hiểu và thông cảm cho việc em đang vô cùng ức chế và khó chịu với những trải nghiệm tồi tệ đã xảy ra trong cái trò chơi đó\dots\ nhưng mà em cũng phải cố gắng hít sâu bình tĩnh lại bớt đi chứ, bạo lực không giải quyết được gì đâu.''

Còn Marine thì ôm đầu thu mình lại sợ hãi tột độ, thái độ khép nép này hoàn toàn trái ngược một trăm tám mươi độ so với cái tính cách mụ giám sát viên hắc ám, hống hách mà cô nàng vừa mới trưng ra trong thế giới trò chơi điện tử ban nãy: ``Ơ kìa\dots\ Em gái ơi, em bớt nóng hít sâu bình tĩnh lại đi, có được không? Bọn chị có làm gì đắc tội với em đâu mà em nổi điên đáng sợ thế!?''

Ryuuhou nghe tiếng cất lên liền khựng tay lại, quay phắt người sang lườm Marine với một ánh mắt hình viên đạn như thể muốn ``ăn tươi nuốt sống'' đối phương. Rõ ràng là con bé vẫn còn ghim thù hậm hực rất sâu đối với cái nhân vật giám sát viên hắc ám do chính cô nàng này vào vai: ``Chị thì không làm gì! Nhưng cái con mụ nhân vật giám sát viên cẩu huyết của chị trong cái game rác đó thì lại cực kỳ xứng đáng bị ăn gậy đập nát sọ!''

Nói là làm, con bé lập tức vung cao cây gậy bóng chày lên không trung, suýt chút nữa là đã không kìm được cơn giận mà giáng thẳng một đòn xuống đầu Marine nếu không nhờ có sự can thiệp ứng cứu kịp thời của Subaru và Mel lao ra cản lại.

Nàng Vịt hốt hoảng chạy tới dang hai tay ra, dũng cảm dùng thân mình làm lá chắn bảo vệ cho đàn em, hét lớn: ``G-Gượm đã! Dừng tay lại ngay lập tức! Em ấy chỉ là bị hệ thống ép buộc vào vai nhân vật giả tưởng thôi mà! Oan có đầu nợ có chủ chứ!''

Nàng Dơi thì vội vàng ôm chặt lấy cánh tay đang cầm gậy của con bé đang nóng máu kia kéo lại: ``Đúng rồi đó, xin em hãy bình tĩnh đi mà. Đó rốt cuộc cũng chỉ là lỗi lập trình của game thôi, đâu phải là lỗi cố ý của cô ấy ở ngoài đời thực đâu.''

Từ phía đằng sau lưng, Kuon khẽ thở dài một hơi não nề, nãy giờ cậu vẫn khoanh tay đứng lặng lẽ quan sát mọi màn hỗn chiến ẩu đả từ xa: ``Thôi được rồi Ryuuhou, em làm ơn đừng để cơn tức giận làm mất kiểm soát lý trí nữa. Game thì em cũng đã phá đảo thành công lôi người ra được rồi, cái máy console nguyên phạm thì cũng đã bị em đập cho nát bươm, biến dạng đến độ đem đi bán cho mấy bà đồng nát người ta cũng chẳng thèm thu mua rồi. Em tính đập phá luôn cả cái văn phòng này mới chịu vừa lòng sao?''

Nghe những lời phân tích đầy lý lẽ và sự khuyên can nhiệt tình của cả ba người, Ryuuhou cuối cùng cũng chịu hít một hơi thật sâu để điều hòa lại nhịp thở. Con bé từ từ hạ tay xuống, vứt phịch cây gậy bóng chày móp méo xuống mặt sàn, rồi khẽ gật đầu gượng gạo với mọi người: ``\dots Mọi người nói cũng có lý. Được rồi, em sẽ nể mặt mọi người mà tha cho\dots\ Nhưng em xin thề độc với trời đất, từ nay về sau bà mày sẽ éo bao giờ thèm đụng tay vào bất kỳ một cái tựa game rác rưởi nào của cái hãng phát triển này thêm một lần nào nữa!''

Thấy rằng sát thủ cầm gậy cuối cùng cũng đã chịu nguôi ngoai cơn giận và lấy lại được sự bình tĩnh, Marine mới dám lí nhí cất lời cảm ơn, cả cơ thể cô nàng lúc này vẫn còn đang run bần bật vì dư chấn sợ hãi: ``C-Cảm tạ em vì đã khoan hồng không vác gậy đánh chị. Nhưng mà\dots\ chị nghĩ tốt nhất là bây giờ chị nên chuồn lẹ đi chỗ khác cho an toàn đây.''

Nói rồi, cô nàng Hải tặc vội vã co giò chạy thục mạng rời khỏi văn phòng như có ma đuổi. Cả căn phòng họp rộng lớn bỗng chốc rơi vào một khoảng không im lặng đến đáng sợ, xung quanh giờ đây chỉ còn văng vẳng lại những tiếng thở dốc nặng nề mệt mỏi của Ryuuhou và Kuon, cùng với những ánh mắt thương xót, đầy tội nghiệp hướng về phía đống nhựa và mạch điện vỡ nát tươm, thứ mà cách đây vài giờ trước vẫn còn là một chiếc máy chơi điện tử hoàn lặn, nguyên vẹn.

\ct

Một khoảng thời gian ngắn sau khi mọi sự hỗn loạn và ẩu đả đã thực sự tạm lắng xuống bề chìm, Pekora bỗng nhiên khẽ ú ớ vài tiếng trong cổ họng rồi bừng tỉnh dậy. Cô nàng Thỏ ngơ ngác đảo đôi mắt nhìn quanh quất một lượt khung cảnh hoang tàn của văn phòng với vẻ mặt vô cùng khó hiểu: ``\dots À-rế? Rốt cuộc là đã có chuyện gì kinh thiên động địa vừa mới xảy ra ở đây vậy, peko?''

Nhìn kỹ lại và thấy rằng con ngươi trên đôi mắt của nàng Thỏ đã phục hồi lại hình dáng thỏ bình thường, không còn lờ đờ vô hồn nữa, Kuon mới khẽ thở phào nhẹ nhõm một hơi dài, giải thích vô cùng ngắn gọn và súc tích: ``Chúc mừng em nha Peko-chan, em vừa mới trải qua một cơn tái phát trầm trọng rơi vào trạng thái nghiện game u mê hệt như cái bữa trước đấy.''

Nghe thấy lời thông báo chấn động của vị Tổng quản lý, Pekora đứng ngớ người ra kinh ngạc. Cô nàng khựng lại mất vài giây, đưa tay lên trán cố gắng vắt óc lục lọi lại mọi mảnh vỡ ký ức để chắp vá lại câu chuyện. Cuối cùng, cô bỗng bất giác sực tỉnh, đập mạnh tay một cái 'đét' vào đầu như thể vừa mới nhớ ra một điều gì đó vô cùng quan trọng: ``À đúng rồi em nhớ ra rồi, peko! Sáng ngày hôm nay trong lúc dọn dẹp em có tình cờ tìm thấy cái tàn tích máy chơi game cũ mà Rushia-chan đã lỡ tay đập nát nhừ tử từ hôm nọ. Em cứ đinh ninh nghĩ là nó vẫn còn có thể sửa chữa để tái chế xài tạm được, nên em mới tò mò cắm điện thử ngồi nghịch bấm bấm vài cái xem sao\dots\ Ai mà có ngờ được hậu quả lại ra nông nỗi này đâu\dots\ peko\dots''

Hoàn toàn không cần phải kiên nhẫn chờ đợi cho cô nàng Thỏ thốt nốt câu chuyện phân trần của mình, cậu Tổng đã lập tức lạnh lùng nối tiếp lời chặn họng: ``\dots và kết cục nhãn tiền là em tự biến chính bản thân mình thành một con nghiện game dặt dẹo mất trí nhớ, để rồi kéo theo bọn anh phải huy động cả đội ngũ cứu hộ chạy đôn chạy đáo vào sinh ra tử trong game để giải quyết hậu quả do em gây ra đấy.''

Nhận thấy ánh mắt nghiêm nghị và đầy tính sát thương cao từ cậu, nàng Thỏ xanh đành rụt cổ cúi gằm mặt xuống đất, hối lỗi lí nhí thanh minh: ``Ơ\dots\ dạ em biết lỗi rồi, em thành thật xin lỗi mọi người nhiều lắm, peko\dots\ Em thực sự không hề có cố ý gây chuyện tày đình thế này đâu mà\dots''

Cậu Tổng thở dài ngao ngán, nhẹ giọng nói tiếp: ``Dù sao thì cái đống phế liệu đó nay cũng đã bị Ryuuhou đập nát bươm không còn hình thù gì nữa rồi, đợi vài ngày nữa dư dả kinh phí thì anh sẽ xuất quỹ công ty sắm lại cho một cái máy khác đời mới xịn sò hơn. Nhưng mà chốt lại một câu ấy, từ nay về sau em vui lòng tém tém lại giúp anh, đừng có mà cắm đầu cắm cổ cày game quá đà quên ăn quên ngủ nữa nhé.''

Pekora nhận được sự khoan hồng chỉ biết gật đầu lia lịa hệt như một cỗ máy, ngoan ngoãn tuyệt đối không dám mở miệng cãi bướng thêm lấy nửa lời. 

Kuon đưa mắt nhìn lướt qua tất cả những gương mặt tàn tạ, mệt mỏi của mọi người có mặt trong phòng, rồi nghiêm giọng dõng dạc tuyên bố một sắc lệnh mới: ``Nhân tiện trải qua cái vụ tai nạn hy hữu này thì anh cũng xin được thông báo ban hành luật mới luôn. Kể từ giây phút này trở đi, anh sẽ chính thức bổ sung thêm một nội quy nghiêm ngặt mới vào sổ tay công ty: giới hạn thời lượng chơi game giải trí trên tất cả các thiết bị máy chơi game cầm tay hay console tại văn phòng này. Anh hoàn toàn không cần quan tâm là mấy đứa sẽ phải mất bao nhiêu thời gian cày cuốc để chuẩn bị content phục vụ cho việc livestream cá nhân, nhưng nội quy chết ở văn phòng này là: không một ai được phép chơi game liên tục vượt quá mức hai tiếng đồng hồ một ngày. Tất cả đã nghe rõ quy định chưa?''

Mọi người có mặt đều ngoan ngoãn gật đầu cái rụp đồng tình không một ý kiến phản bác. Subaru và Mel thậm chí còn không quên khẽ liếc xéo đôi mắt sang trừng trị Pekora, ngầm phóng ra thông điệp vô hình như muốn mắng vốn: ``Đấy, em thấy chưa, em đúng là bà chúa chuyên gia gây ra đủ mọi thứ rắc rối phiền toái trên đời thật đấy.''

\dots

\dots

\dots

Mặc dù nhóm ba người anh hùng quả cảm cuối cùng cũng đã đồng tâm hiệp lực đánh bại được ả trùm cuối Choco và phá đảo thành công trò chơi. Thế nhưng, bằng một cách thần kỳ hay lỗi hệ thống nào đó, bản thân ả phản diện Choco vẫn bị mắc kẹt vĩnh viễn ở lại bên trong cái thế giới game hỗn loạn mà Rushia đã từng vô tình nhúng tay tạo ra khi sửa máy.

Nàng Y tá đáng thương đưa mắt nhìn quanh quất một lượt khung cảnh hoang tàn, đổ nát của trụ sở tập đoàn công ty sữa lắc (có vẻ như ý thức của cô nàng xui xẻo này vẫn chưa thể thoát khỏi được vai diễn nữ phản diện), bất lực gào lên xé toạc cả không gian tĩnh lặng: ``\textbf{Cái quái quỷ gì đang diễn ra thế này!? Tại sao cả thế giới rộng lớn này lại đột nhiên chỉ còn trơ trọi sót lại mỗi mình ta bị kẹt ở lại đây vậy!? Lũ anh hùng khốn khiếp kia đâu hết rồi!?}''

Cô ả hoảng loạn vung tay thử thi triển đủ mọi loại phép thuật để triệu hồi đàn em, thậm chí ả còn tuyệt vọng đến mức cố gắng hack và can thiệp sâu vào những dòng mã lập trình code của thế giới trò chơi để tìm đường thoát thân, thế nhưng mọi nỗ lực đều như muối bỏ bể, hoàn toàn không nhận được bất kỳ một tín hiệu phản hồi nào từ máy chủ.

Giữa lúc đang chìm trong sự hoảng loạn tột cùng, bỗng nhiên, một giọng nói điện tử ồm ồm lạnh lẽo phát ra từ hệ thống quản lý trò chơi vang vọng khắp không gian, làm cho nữ Y tá giật bắn mình thon thót.

\begin{center}
  ``Kính chào mừng người chơi đã bước vào chế độ cày cuốc mới mang tên: `Thử thách sinh tồn khắc nghiệt chốn công sở!' Nhiệm vụ bắt buộc của bạn là: Phải tham gia lao động sản xuất liên tục trong suốt 900 giờ đồng hồ để thu thập đủ điểm tích lũy và mở khóa lối thoát khỏi không gian này.''
\end{center}

Ngay lập tức, hiện lên lù lù ngay trước mắt cô ả là một bảng đồng hồ đếm ngược với con số {\digi 900} khổng lồ chói lóa được tính theo đơn vị giờ đồng hồ, nối tiếp theo sau đó là vô vàn những con số 0 tròn trĩnh đếm lùi nhảy múa liên tục từ phút, giây và phần trăm tích tắc.

Choco nhìn bảng đếm ngược mà thảng thốt rụng rời chân tay: ``9\dots\ 900 giờ lao động liên tục cơ á!? Các người bị điên hết cả rồi sao! Ta đường đường là một nữ giám đốc điều hành chủ tập đoàn công ty vĩ đại cơ mà! Làm gì có cái lý lẽ nào bắt một người làm chủ bệ vệ như ta lại phải hạ mình đi cày cuốc làm ba cái việc lao động chân tay thấp hèn như thế này được!''

Thế nhưng, như để cố tình chọc tức và trêu ngươi cô nàng phản diện kiêu ngạo vẫn còn đang bị mắc kẹt vô vọng ở nơi khỉ ho cò gáy này, hệ thống thông báo lại tiếp tục lạnh lùng phát ra một dòng chữ tàn nhẫn:

\begin{center}
  ``Lưu ý cực kỳ quan trọng: Cứ mỗi khi hệ thống ghi nhận người chơi thốt ra lời phàn nàn chửi bới, thời gian lao động bắt buộc sẽ tự động bị phạt cộng dồn tăng thêm 100 giờ đồng hồ.''
\end{center}

Choco vừa nghe xong cái hình thức xử phạt đến mức nực cười và vô lý chưa từng có trong lịch sử ngành game đó, lập tức đánh mất luôn phong thái điềm tĩnh, hét lớn lên đầy tức giận điên cuồng: ``\textbf{Không bao giờ có cái mùa xuân đó đâu nhé! Lũ khốn nạn các ngươi không có quyền được phép đối xử với một vị chủ tịch như ta một cách vô nhân đạo như vậy! Đánh chết ta cũng méo bao giờ thèm làm đâu nghe chưa!}''

*BÍP* Một âm thanh thông báo vang lên khô khốc, cắt ngang tiếng chửi rủa.

\begin{center}
  ``Hệ thống ghi nhận vi phạm. Thời gian lao động hình phạt đã được tự động cộng dồn và cập nhật tăng lên thành 1000 giờ.''
\end{center}

Ngay lập tức, con số {\digi 900} đỏ rực hiển thị ở cột mốc giờ liền tự động nhảy vọt lên thành con số {\digi \ 1000} đầy ám ảnh và tuyệt vọng.

Nàng Y tá quyền lực lúc này hoàn toàn bất lực, đành suy sụp ngồi bệt thụp xuống mặt nền đất, hai tay ôm đầu vò rối tung cả mái tóc vàng óng ả. Trong tông giọng nghẹn ngào chứa đầy sự tức tối và bất lực đến cùng cực, cô ả chỉ còn biết thốt ra đúng một câu chửi rủa trong nước mắt: ``Ta thực sự căm thù cái thế giới chết tiệt này\dots\ Lũ rác rưởi các ngươi không có quyền\dots\ làm thế này với ta cơ mà\dots!''

Ở bên ngoài thế giới trò chơi thực tại, cả nhóm sáu người đứng vây quanh màn hình tivi, cùng nhau nhìn chăm chú vào cái khung cảnh thảm thương mà ả trùm cuối kiêu ngạo đang phải hứng chịu. Mỗi người trong số họ đều mang trong mình một tâm trạng và dòng suy nghĩ khác nhau.

Mel tỏ vẻ thương hại, mủi lòng xót xa cho tình cảnh của Choco: ``Ôi, nhìn tội nghiệp cho Choco-sensei ghê á. Giá như mà ngay từ đầu cô ấy không bị hệ thống ép buộc phải phân vào đóng cái vai ác nhân thì kết cục đã tốt biết mấy\dots''

Subaru thì ngược lại, khoanh tay đứng nhìn với vẻ mặt vô cùng hả hê và thỏa mãn với hình phạt thích đáng này: ``Thôi đi Mel-senpai ơi, chị hoàn toàn không cần phải tỏ ra thương xót từ bi với ả ta làm gì đâu. Cái loại tư bản bóc lột đó cực kỳ xứng đáng bị chịu cái hình phạt đày đọa này để cho biết mùi đời!''

Marine, người cũng vừa mới vội vã quay trở lại văn phòng sau một hồi chạy ra ngoài hít thở khí trời lấy lại bình tĩnh, cũng lập tức vung nắm đấm thốt lên đầy nhiệt tình hưởng ứng: ``Đúng đó! Cho cô ta nếm trải mùi vị bị bóc lột cho biết thế nào là lễ độ đi! Cũng chỉ tại ả ta là đầu sỏ thiết kế ra cái kịch bản quái thai đó mà tôi vừa mới suýt chút nữa bị cái cô em gái bạo lực nhà anh vụt gậy đánh cho vỡ đầu đấy!''

Pekora đứng ở ngoài rìa khẽ nghiêng đầu, đôi mắt tỏ vẻ vô cùng khó hiểu trước thái độ thù hằn gay gắt của cả ba người đồng nghiệp. Thế nhưng, khi nhìn vào những cảnh tượng đày đọa khổ sai mà Choco đang phải cắn răng gánh chịu trên màn hình, cô nàng Thỏ cũng không kìm được mà khẽ nhếch mép cười tủm tỉm: ``\dots Thôi thì kệ đi, tự dưng coi mấy cái cảnh này cũng mang tính chất giải trí phết đấy chứ nhỉ, peko.''

Kuon và Ryuuhou đứng khoanh tay dựa lưng vào tường, dán mắt nhìn màn hình trò chơi, cả hai anh em không hẹn mà gặp cùng khẽ nhếch mép cười khẩy một cách đầy mãn nguyện hệt như hai vị thần phán xét.

Cậu Tổng quản lý nhếch mép buông lời: ``Đấy em thấy chưa, có vẻ như nãy anh cản em lại không đập nát hoàn toàn bảng mạch hệ thống lại là một quyết định vô cùng sáng suốt đấy. Cứ để mặc ả phản diện đó ở lại trong đấy cày cuốc sml để tự rút ra bài học đắt giá cho cái thói kiêu ngạo của mình đi.''

Cô em gái game thủ gật gù tán thành, cười nhạt bồi thêm: ``Ừ, công nhận. Có khi cái việc giam cầm đày đọa lao động không lối thoát này lại là một cách trừng phạt nhân đạo và công bằng nhất dành cho lũ tư bản bóc lột mỏ hỗn đấy.''

Tất cả mọi người trong phòng đều đồng loạt hướng ánh mắt nhìn cô nàng Y tá kiêu ngạo nay vẫn chưa thể thoát khỏi vai diễn đang vừa phải gào thét kêu ca phàn nàn vừa phải cắn răng cam chịu làm cái công việc lao động khổ sai chân tay, trên gương mặt ai nấy đều toát lên một vẻ khoan khoái như thể toàn bộ những cục tức, những sự bực dọc dồn nén trong người từ nãy đến giờ đều đã được đấm tan biến giải thoát sạch sẽ. 

Cậu Tổng quản lý Kuon tuy ngoài mặt thì cười cười tỏ vẻ tàn nhẫn vậy thôi, nhưng sâu thẳm trong lòng vẫn âm thầm tự nhẩm tính toán đường lui: ``\textit{Chửi thì cứ chửi cho sướng miệng vậy thôi, chứ dù sao đi nữa một lát nữa chắc chắn mình cũng vẫn phải lén lút tìm cách viết mã code để lôi cổ em ấy ra khỏi cái thế giới ảo đó mới được\dots\ Chứ cứ để nhốt em ấy ở trong cái vòng lặp 1000 giờ điên rồ đó không khéo em ấy sang chấn tâm lý rồi ngoẹo luôn trong đó chứ chẳng đùa đâu\dots}''

\end{document}